\chapter{Clause Types and Modality}\label{ch:clause-types}

\section{\ird{By}: the canonical exponent of \texorpdfstring{\synhead{Pol\textsuperscript{0}}}{Pol0}}
\label{sec:by-pol-head}

Interrogative \ird{by} is analysed as the overt exponent of
\synhead{Pol\textsuperscript{0}}, below the ordinary placement edge
\synhead{OrdP} and the independently licensed \synhead{FocP} and
\synhead{TopP} (§\,\ref{sec:th-clause-structure}). Its operator types the
clause as interrogative and takes scope over the comment domain, whereas
additive \ird{oh} and affirmative \ird{mó} contribute meanings within
the comment. Their linear proximity therefore coexists with distinct
structural positions and scope relations. Pol\textsuperscript{0} is
locally head-initial, yielding \ird{by} before its complement; ordinary
arguments and modifiers move to \synhead{OrdP} above the particle field
and are pronounced before it
(§\,\ref{sec:tc-adverbials}).

Prosodically, \ird{by} is a weak independent particle: its monosyllable may
be deaccented in an ordinary question or made prominent for insistence.
Neither its position after ordinary constituents nor its scope establishes
obligatory attachment to a preceding nominal or to following material. Its
phrasing does not alter
the lexical stress or subsidiary footing of neighbouring words
(§§\,\ref{sec:stress} and~\ref{sec:secondary-stress}). An ordinary polar
question has a low nuclear target on the prominent stressed
syllable followed by a final rise, represented by L* H\%. The rise may
extend through post-stressed syllables to the phrase edge, while a marked
falling contour remains compatible with overt interrogative \ird{by}
(§\,\ref{sec:intonation}). The same rising contour is available with
null Pol\textsuperscript{0}; its melody expresses the interrogative use
without supplying a second clause-typing mechanism.

The analysis of \ird{by} as \synhead{Pol\textsuperscript{0}} entails the following
structural claims. First, \ird{by} is the overt exponent of \synhead{Pol\textsuperscript{0}}
in matrix interrogatives generally: it types the whole clause as interrogative
and scopes over the full \synhead{IP}, including any comment-domain material. Second,
\synhead{Pol\textsuperscript{0}} may be phonologically null. In polar questions, overt
\ird{by} is the ordinary and unmarked explicit strategy; in content questions,
both overt \ird{by} and null \synhead{Pol\textsuperscript{0}} are standardly available
(see §\,\ref{sec:th-wh-interrogation}). The choice between them in content
questions reflects a difference in discourse profile --- how explicitly and
directly the clause is presented as a question in its interactional context ---
rather than a grammatical asymmetry. Third, \synhead{Pol\textsuperscript{0}} is not a
focus-licensing head; Spec,\synhead{FocP}, above \synhead{OrdP}, hosts
a marked fronted wh-phrase or another focused constituent on independent
grounds, as discussed in
§\,\ref{sec:th-focus}. The movement of a wh-phrase to the left-peripheral
focus position is independent of whether \synhead{Pol\textsuperscript{0}} is overt or
null. Ordinary wh-placement in \synhead{OrdP} also precedes \ird{by},
without constituting dedicated wh-fronting; interrogative interpretation is
connected to the independently licensed lower occurrence.

The contribution of the Class~1 phasal particles \ird{vy} and \ird{so} to a
polar question is therefore independent of \ird{by} itself. Where \ird{vy}
and \ird{so} specify the phasal profile of the predicate — whether an event
is already attained or an interval still open — \ird{by} contributes the
clause-level interrogative force. \synhead{PolP} dominates the Comm-field
and \synhead{PhasP}, while outer \synhead{EvidP} and
\synhead{SpeechActP} stand above it: a mirative question
(\ird{oče~by}) presents the polar question as counterexpectational; a
reportative question (\ird{či~by}) presents it as relayed from another
source. Their high position is consistent with those interpretations, but
does not by itself distinguish surprise at an answer from surprise at the
asking of a question, or attribution of a proposition from attribution of
the question itself. Nor does it determine whose evidence or uncertainty
is relevant. The established high scope therefore leaves these finer
interpretive relations unspecified; the distinction between structural
height and the object of an attitude is developed in
§\,\ref{sec:th-polar-questions}.

\pex
\a
\begingl
\gla Marek by zdrávime?//
\glb Marek.\Dir{} \textsc{q} sleep-\Prog{}//
\glft \trsl{Is Marek sleeping?}//
\endgl
\a
\begingl
\gla Janek by vy pjášček?//
\glb Janek.\Dir{} \textsc{q} already eat-\Antip{}-\Pf{}//
\glft \trsl{Has Janek already eaten?}//
\endgl

\xe

\noindent In ordinary conversation, intonation and wh-marking can make the
interrogative interpretation recoverable when Pol\textsuperscript{0} is
phonologically null, but overt \ird{by} remains the clearest
and most explicit matrix interrogative marker across both polar and content
questions. Its use is especially expected in careful, explicit, or formally
written style, in interactionally foregrounded questions, and whenever the
speaker wants the interrogative framing to be unambiguous to the addressee.
The full structural analysis of polar and content questions --- including the
four-combination paradigm of \ird{by}/null Pol crossed with
marked wh-fronting/ordinary wh-placement --- is in §\,\ref{sec:th-polar-questions} and
§\,\ref{sec:th-wh-interrogation}.

\subsection{Polar questions and higher particles}
\label{sec:th-polar-questions}

Polar questions ordinarily employ overt \ird{by}, whose clause-typing
contribution remains independent of a phasal operator:

\pex
\a
\begingl
\gla Janek by vy ščének?//
\glb Janek.\Dir{} \textsc{q} already arrive-\Pf{}//
\glft \trsl{Has Janek already arrived?}//
\endgl
\xe

Zone~\textsc{i} particles follow ordinary constituents in questions,
with \ird{by} typing the lower clause as interrogative. Mirative \ird{oče}
supplies the surprise or counterexpectational framing illustrated below,
whereas reportative \ird{či} supplies attributed or relayed framing. Their
positions above \synhead{PolP} represent that general relation, but do not
independently determine whether surprise concerns a possible answer, an
earlier expectation, or the occasion for asking, nor whether a reportative
question relays someone else's question or seeks an answer supported by a
report. Their placement establishes the high framing without
specifying the distribution of all these more particular readings.

\pex
\a
\begingl
\gla Tak tóm Janku oče by vy koupnik?//
\glb this book.\Dir{} Janek.\Trs{} \textsc{mir} \textsc{q} already buy-\Trs{}-\Pf{}//
\glft \trsl{Has Janek, surprisingly, already bought this book?}//
\endgl
\xe

The translation illustrates surprise associated with Janek's possible
purchase, while leaving the purchase itself at issue; it does not assert that
he bought the book. No further reading follows simply from the relative
pronunciation of \ird{oče} and \ird{by}. This distinction matters in the
Tagalog comparison: AnderBois analyses interrogative \textit{pala} as a
revelation about performing the question update, including remembering an
intention to ask or recognising the question's relevance
\cite[§5.2]{anderbois2023pala}, while Tan Almazán distinguishes reported
questions from questions anticipating reportative evidence for an answer
with \textit{daw} \cite[§4.3.4]{tanalmazan2019}. These findings identify
different interpretive dependencies to be distinguished; they do not assign
the Tagalog readings to Iridian \ird{oče} or \ird{či}.

A constituent may be fronted to Spec,\synhead{FocP} or set as a topic in
Spec,\synhead{TopP} inside a polar question, while \ird{by} in
Pol\textsuperscript{0} continues to type the clause as interrogative. The
fronting of a non-wh constituent inside a polar question is not a special
interrogative mechanism but ordinary information structure operating inside an
interrogative clause. A direct-case topic may remain in Spec,\synhead{TopP}; a
narrow-focus non-wh constituent may move to Spec,\synhead{FocP}; either
operation is independent of the interrogative clause typing supplied by
\ird{by}, whose operator supplies interrogative force independently of the
position chosen for the argument or adjunct.


\section{Content questions and \textit{wh}-interrogation}
\label{sec:th-wh-interrogation}

\subsection{Licensing of content questions: \synhead{PolP} and wh-operators}
\label{sec:th-single-wh}

Content questions combine interrogative clause typing with a wh-expression
whose lower occurrence supplies a variable in the same finite domain. In this
analysis, interrogative Pol binds that variable; neither ordinary placement
nor dedicated wh-fronting supplies interrogative force by itself.
\textsc{Pol}\textsuperscript{0} may be phonologically null in neutral content
questions, while remaining the clause-typing head.

In neutral content questions, wh-expressions undergo ordinary placement
 with the other arguments and modifiers, without dedicated wh-fronting.
 Their lower chain occurrences supply interrogative variables within the
 scope of Pol. A single wh-expression may additionally move to
 Spec,\synhead{FocP}, at the left edge of the ordinary placement field,
 when singled out against a salient frame or assigned contrast, insistence
 or other discourse prominence. At most one marked wh-fronting dependency
 is licensed per clause; ordinary placement of several wh-expressions does
 not constitute multiple marked frontings.

\pex
\a
\begingl
\gla Marek jena mulša?//
\glb Marek.\Dir{} where live-\Cont{}//
\glft \trsl{Where does Marek live?} \lingcontext{neutral; without dedicated wh-fronting}//
\endgl
\a
\begingl
\gla Jena Marek mulša?//
\glb where Marek.\Dir{} live-\Cont{}//
\glft \trsl{Where does Marek live?} \lingcontext{marked; wh to \synhead{SpecFocP} for discourse prominence}//
\endgl
\a
\begingl
\gla Tak tóm jeda koupnik?//
\glb this book.\Dir{} who buy-\Pf{}//
\glft \trsl{Who bought this book?} \lingcontext{neutral; without dedicated wh-fronting}//
\endgl
\a
\begingl
\gla Jeda tak tóm koupnik?//
\glb who this book.\Dir{} buy-\Pf{}//
\glft \trsl{Who bought this book?} \lingcontext{marked; wh to SpecFocP}//
\endgl
\xe

A fronted wh-phrase does not itself license interrogative force; its movement
signals discourse prominence, contrast, or singled-out status of the
wh-variable. The interrogative clause type is established by the \synhead{PolP}
layer independently of whether any fronting has occurred. A content question is
therefore structurally defined by the conjunction of an interrogative operator
in \synhead{PolP} and a wh-expression interpreted within its scope, not by the position
of the wh-phrase at the surface.

Matrix content questions permit both overt and null Pol\textsuperscript{0}, and
these are \emph{both} standard options, not a default with an emphatic variant.
The two differ in discourse profile. A content question with null
Pol\textsuperscript{0} is the structurally lighter form, appropriately neutral
in ordinary information-seeking contexts. A content question with overt \ird{by}
is more explicitly typed as a direct question: it is more interactionally
foregrounded, conversationally engaged, or oriented toward a specific addressee
in the current exchange. The difference is a pragmatic tendency, not a semantic
or grammatical opposition. Both forms are equally grammatical in standard
Iridian.

Content questions permit a concluding fall, with nuclear prominence
determined by the focus domain rather than by an obligatory question rise
(§\,\ref{sec:intonation}). A prominent wh-expression retains its initial
lexical stress whether it receives ordinary placement or marked fronting,
and overt \ird{by} does not impose a different nuclear contour.

Both relative orders of Marek and the wh-phrase occur with null and overt
Pol\textsuperscript{0}. The resulting four combinations distinguish the
explicitness of clause typing from the discourse prominence of the wh-phrase.
The paradigm distinguishes the two choices in the surface patterns.

The following paradigm distinguishes ordinary placement from marked
wh-fronting, with both overt and null Pol:

\pex\label{ex:wh-four-combinations}
\a
\begingl
\gla Marek jena mulša?//
\glb Marek.\Dir{} where live-\Cont{}//
\glft \trsl{Where does Marek live?} \lingcontext{null Pol; without dedicated wh-fronting: neutral}//
\endgl
\a
\begingl
\gla Marek jena by mulša?//
\glb Marek.\Dir{} where \textsc{q} live-\Cont{}//
\glft \trsl{Where does Marek live?} \lingcontext{overt \ird{by}; without dedicated wh-fronting: explicit, interactionally foregrounded}//
\endgl
\a
\begingl
\gla Jena Marek mulša?//
\glb where Marek.\Dir{} live-\Cont{}//
\glft \trsl{Where does Marek live?} \lingcontext{null Pol; wh fronted: wh-phrase discourse-prominent}//
\endgl
\a
\begingl
\gla Jena Marek by mulša?//
\glb where Marek.\Dir{} \textsc{q} live-\Cont{}//
\glft \trsl{Where does Marek live?} \lingcontext{overt \ird{by}; wh fronted: both clause and wh-phrase foregrounded}//
\endgl
\xe

The same four-way system extends to transitive and oblique content
questions. The wh-word type (agent, patient, location, instrument, and so on)
does not restrict the availability of any combination; all four remain possible
across the full inventory of wh-operators in Table~\ref{tab:nm-interrogatives}.

Zone~\textsc{i} particles follow ordinary and marked wh-expressions,
 although their operator positions remain above Pol:

\pex
\a
\begingl
\gla Marek jena hladě mulša?//
\glb Marek.\Dir{} where \textsc{infr} live-\Cont{}//
\glft \trsl{Where would Marek presumably be living?}\\
      \lingcontext{\ird{hladě} in \synhead{EvidP} above \synhead{PolP}; without dedicated wh-fronting}//
\endgl
\xe

The inferential framing leaves Marek's location unresolved. Neither the
position of \ird{hladě} nor its English paraphrase independently
specifies whose evidence would support an answer; that dependency remains
distinct from the particle's established upper-peripheral placement, as in
the polar-question discussion in §\,\ref{sec:th-polar-questions}.

\subsection{Multiple \textit{wh} and ordinary placement}
\label{sec:th-multiple-wh}

When the clause contains more than one wh-phrase, the ordinary pattern is for
all wh-phrases to undergo ordinary placement without dedicated wh-fronting.
At most one wh-phrase may undergo marked fronting,
and only when that phrase is assigned specially marked focus
within the pair-list interpretation. The remaining wh-phrases are interpreted as
interrogative variables through their lower chain occurrences, licensed by the clause-level
interrogative operator (overt \ird{by} or a null \synhead{PolP} head).

\pex
\a
\begingl
\gla Jede ježe koupnik?//
\glb who what buy-\Pf{}//
\glft \trsl{Who bought what?}//
\endgl
\a
\begingl
\gla Jede jena mulša?//
\glb who where live-\Cont{}//
\glft \trsl{Who lives where?}//
\endgl
\xe

\noindent In these multiple-wh clauses neither phrase undergoes marked
wh-fronting; their order does not by itself establish their exact positions
within the lower clause. Their lower occurrences are bound by the null interrogative
operator in \synhead{PolP}. The licensing principle relates a wh-expression's eligible lower occurrence
within its own finite domain to interrogative Pol; ordinary or marked
pronunciation above Pol does not sever that relation. It supplies no
independent licence for a dependency across a finite or nominal boundary.
Overt fronting marks the phrase as the salient or
foregrounded variable in the pair-list; its landing site is Spec,\synhead{FocP}, above \synhead{OrdP}, as
illustrated in \ref{ex:wh-four-combinations}. Fronting does not itself establish the
interrogative type of the clause, which is always a function of \synhead{PolP}.

\subsection{Tag, echo, and indirect questions}
\label{sec:th-tag-echo}

For indirect questions embedded under \ird{poby}, see §\,\ref{sec:poby}. For
embedded content questions in \ird{poby}-complements, the same distributional
principle applies as in matrix clauses: the embedded wh-phrase may undergo ordinary placement
(the neutral option) or additional marked movement to the embedded focus edge
for marked focus. Wh-fronting within the embedded clause's own periphery is
therefore not obligatory; it signals the same kind of discourse prominence as
fronting in matrix questions. The \ird{poby} complementizer marks the embedded
clause as a non-assertive complement regardless of whether the internal
wh-phrase has fronted. The full analysis of \ird{poby} and its complement types
is in §\,\ref{sec:poby}.


\section{The directive domain}\label{sec:mood-directive}

The imperative, the prohibitive, and the hortative form a unified directive
domain: a compact finite subsystem used for commands, negative commands,
suggestions, proposals, soft requests, permissive uses, and wishes. Where the
indicative provides the full five-way aspect system and the split-ergative
alignment described in Section~\ref{sec:aspect}, directive moods carry a reduced
two-way aspect-domain contrast in Asp.Align: the imperfective domain
($\varnothing$ exponent) and the perfective domain
(\ird{-ov-} exponent), both subject to directive orientation. The directive mood suffix follows
reduced Asp.Align, outside any applicative or causative morphology
(§\,\ref{sec:vm-mood-asp}).

Directive Mood selects the theme as the direct-case participant in both
aspectual domains, while an agent or causer with a distinct direct participant
takes transitive case. With no theme, an applied participant or causee can be
selected; an underived intransitive retains its sole argument in direct case.
The understood subject of a command is normally recoverable from the speech
situation, although its omission does not alter this case frame. Applied
participants and causees must be overt in directives; when a theme is retained,
they take indirect case. Antipassive morphology is excluded by a Mood--Voice
licensing restriction (§\,\ref{sec:mood-directive-valency}).

\subsection{Imperative and prohibitive}\label{sec:mood-imp}

The imperative (\ird{-ím}) is the ordinary mood for direct commands,
instructions, warnings, and immediate discourse-management directives addressed
to the hearer. The prohibitive (\ird{-é}) is its dedicated negative counterpart,
used for negative commands and preventative directives. The two are treated
together because they share the same null-subject pattern and the same case
geometry.

The subject of both the imperative and the prohibitive is second person and is
syntactically understood but ordinarily unexpressed; no overt pronoun is
required in fluent discourse. The theme or patient of the clause occupies the
direct case as the clause's single core direct nominal:

\pex
\a
\begingl
\gla Trava pjaštím.//
\glb bread.\Dir{} eat-\Imp{}//
\glft \trsl{Eat the bread.}//
\endgl
\a
\begingl
\gla Tak hrabním.//
\glb \Dem{}.\Prox{}.\Dir{} hear-\Imp{}//
\glft \trsl{Hear this.}//
\endgl
\xe

The prohibitive directs the addressee to refrain from the action and excludes
the indicative negative marker \ird{za-}. It is analysed as an exponent of
negative Mood (§\,\ref{sec:vm-moodp}):

\pex
\a
\begingl
\gla Trava pjašté.//
\glb bread.\Dir{} eat-\Proh{}//
\glft \trsl{Do not eat the bread.}//
\endgl
\xe

Overt second-person subjects are permitted in both the imperative and the
prohibitive, but they are pragmatically marked, expressing foregrounding,
contrast, insistence, or rebuke toward the addressee. The second-person singular
pronoun is \ird{já} in the direct and \ird{jí} in the transitive; the
second-person plural is \ird{tová} in the direct and \ird{tvy} in the
transitive. In a transitive imperative or prohibitive clause where an overt
direct object is already present, an overt speech-act participant subject
appears in the transitive case rather than the direct, so that the affected
participant may remain the clause's sole direct nominal:

\pex
\a
\begingl
\gla Jí trava pjaštím.//
\glb 2\textsc{sg}.\Trs{} bread.\Dir{} eat-\Imp{}//
\glft \trsl{You eat the bread.} \lingcontext{marked; overt 2\textsc{sg} subject in \Trs{}}//
\endgl
\a
\begingl
\gla Tvy trava pjašté.//
\glb 2\textsc{pl}.\Trs{} bread.\Dir{} eat-\Proh{}//
\glft \trsl{You lot do not eat the bread.} \lingcontext{marked; overt 2\textsc{pl} subject in \Trs{}}//
\endgl
\xe

\noindent The agent's transitive case belongs to the directive frame with a
separately selected theme; it is retained when the agent is silent. An
underived intransitive has no distinct theme, so its sole participant takes
direct case.

\subsection{Hortative}\label{sec:mood-hort}

The hortative covers a broader range of directive uses than the imperative:
inclusive proposals of the type \textit{let us}, softened direct commands,
permissive uses of the type \textit{let the child come}, and optative
expressions of wish or hope. It serves as the default polite register in
ordinary speech, typically preferred over the imperative in everyday
conversation. The hortative takes two forms: the positive \ird{-ká} and the
negative \ird{-ku}. Like the imperative and prohibitive, it excludes the
ordinary aspect suffixes and operates within the single-direct principle of the
directive domain.

The hortative defaults to a null subject whose interpretation is supplied by
clause type and context. In proposal uses the null subject is most readily read
as first person plural, inclusive of the addressee; in softened directive uses
the same null subject may receive a second-person reading. Null-subject
hortatives are the unmarked form:

\pex
\a
\begingl
\gla Pjaštká.//
\glb eat-\Hort{}//
\glft \trsl{Let's eat.} \lingcontext{null subject; inclusive proposal}//
\endgl
\a
\begingl
\gla Tak hrabnká.//
\glb \Dem{}.\Prox{}.\Dir{} hear-\Hort{}//
\glft \trsl{Please hear this.} \lingcontext{null subject; softened directive with direct object}//
\endgl
\xe

The negative hortative \ird{-ku} is used for proposals to refrain, cautious
collective restraint, and mild negative requests, with the suffix expressing
negative hortative Mood and excluding the indicative marker \ird{za-}:

\pex
\a
\begingl
\gla Tak pjaštku.//
\glb \Dem{}.\Prox{}.\Dir{} eat-\Neg{}\Hort{}//
\glft \trsl{Let's not eat this.} \lingcontext{negative hortative; collective restraint}//
\endgl
\xe

Overt subjects in the hortative fall into two structurally distinct classes. The
first class comprises first-person plural and second-person subjects in
proposal or softened-command uses. As in the imperative and prohibitive,
these subjects take transitive case when the clause licenses a distinct
participant as direct:

\pex
\a
\begingl
\gla Mně trava pjaštká.//
\glb 1\textsc{pl}.\Trs{} bread.\Dir{} eat-\Hort{}//
\glft \trsl{Let us eat the bread.} \lingcontext{overt 1\textsc{pl} subject in \Trs{}}//
\endgl
\a
\begingl
\gla Jí tak hrabnká.//
\glb 2\textsc{sg}.\Trs{} \Dem{}.\Prox{}.\Dir{} hear-\Hort{}//
\glft \trsl{Please, you hear this.} \lingcontext{marked; overt 2\textsc{sg} subject in \Trs{}}//
\endgl
\xe

The second class comprises overt lexical subjects in optative or permissive
hortatives, where the clause expresses a wish or permission directed toward a
third participant rather than a command to the addressee. Such clauses are
typically intransitive. Their lexical subject occupies the direct case as the
clause's single core direct nominal --- the same position that in transitive
directive clauses would be held by the theme or patient. As the sole argument, it takes
direct case:

\pex
\a
\begingl
\gla Byl ščenká.//
\glb child.\Dir{} arrive-\Hort{}//
\glft \trsl{May the child come.} \lingcontext{intransitive optative; lexical subject in \Dir{}}//
\endgl
\a
\begingl
\gla Je ščenká.//
\glb \Dem{}.\Dist{}.\Dir{} arrive-\Hort{}//
\glft \trsl{May that come to pass.} \lingcontext{intransitive optative; distal demonstrative in \Dir{}}//
\endgl
\xe

\noindent A hortative subject thus receives the case licensed by its argument
frame: transitive with a distinct direct-case theme, direct as the sole
participant of an intransitive. The contrast between commands addressed to
speech-act participants and wishes concerning a third participant does not
itself determine case.

\subsection{Case and subject expression in directive clauses}\label{sec:mood-directive-case}

Directive clauses assign direct case to the theme in both aspectual domains.
An agent or causer takes transitive case when a distinct participant is selected
as direct, whereas the sole argument of an underived intransitive remains
direct. The orientation is imposed by directive Mood and therefore persists
whether the subject is pronounced or omitted. The understood second-person
subject of an imperative or prohibitive is normally silent because the speech
situation identifies it; an overt pronoun can contrast or insist on the
addressee without changing the theme's case:

\pex\label{ex:val-directive-simple}
\a
\begingl
\gla Trava pjaštím.//
\glb bread.\Dir{} eat-\Imp{}//
\glft \trsl{Eat the bread.} \lingcontext{null subject; theme in \Dir{}}//
\endgl
\a
\begingl
\gla Jí trava pjaštím.//
\glb 2\textsc{sg}.\Trs{} bread.\Dir{} eat-\Imp{}//
\glft \trsl{You eat the bread.} \lingcontext{overt subject in \Trs{}; theme in \Dir{}}//
\endgl
\xe

\noindent An overt lexical subject of an intransitive optative or permissive
hortative takes direct case as the clause's sole argument. The transitive
marking of the agent in the second example thus belongs to a clause with a
distinct selected theme; it is independent of grammatical person and of the
choice to pronounce the agent.

\subsection{Valency-changing morphology in directive clauses}\label{sec:mood-directive-valency}

Applicative and causative morphology are compatible with the directive moods.
They precede reduced Asp.Align and the directive mood suffix, so that the
IPFV forms place Mood directly after the derivational suffix, whereas PFV
\ird{-ov-} intervenes before Mood (§\,\ref{sec:vm-mood-asp}).

The applicative (\ird{-éb-}) is licensed in directive clauses only if the
beneficiary is overtly realised. With a retained theme, the beneficiary is an
indirect-case DP, while the theme is direct and any overt agent is transitive.
This Mood-conditioned selection supersedes the simple indicative applicative
frame. A \ird{my} phrase expresses the separate postpositional strategy:

\pex\label{ex:val-directive-appl}
\a
\begingl
\gla Trava bylam pjaštébím.//
\glb bread.\Dir{} child.\Ind{} eat-\Appl{}-\Imp{}//
\glft \trsl{Eat the bread for the child's benefit.} \lingcontext{applicative; beneficiary in \Ind{}}//
\endgl
\a
\begingl
\gla Jí trava bylam pjaštébím.//
\glb 2\textsc{sg}.\Trs{} bread.\Dir{} child.\Ind{} eat-\Appl{}-\Imp{}//
\glft \trsl{You eat the bread for the child's benefit.} \lingcontext{overt subject in \Trs{}}//
\endgl
\xe

\noindent Without an overt beneficiary the applicative is ungrammatical in
directive clauses:

\pex
\rec{trava pjaštébím} \trsl{intended: `Eat bread for someone's benefit.'} \lingcontext{applicative without beneficiary}
\xe

The rhotic causative likewise requires an overtly realised causee in
directive clauses. With a retained lexical theme, the causee takes indirect case
and the theme remains direct; in a causative of an underived intransitive, the
causee itself is direct. An overt causer is transitive in either frame.
Compact \ird{-r-} precedes vowel-initial \ird{-ím} and \ird{-é},
whereas expanded \ird{-er-} precedes consonant-initial hortative
\ird{-ká} and \ird{-ku}: \ird{pjaštérká}, \ird{pjaštérku}.
When PFV \ird{-ov-} intervenes, the compact grade is retained, as in
\ird{pjáštrovím} and \ird{pjáštrovká}. Lexically expanded stems,
including statives in \ird{-n-}, retain \ird{-er-}
(§\,\ref{sec:causative-grades}):

\pex\label{ex:val-directive-caus}
\a
\begingl
\gla Trava bylam pjáštrím.//
\glb bread.\Dir{} child.\Ind{} eat-\Caus{}-\Imp{}//
\glft \trsl{Make the child eat the bread.} \lingcontext{causative; causee in \Ind{}}//
\endgl
\a
\begingl
\gla Jí trava bylam pjáštrím.//
\glb 2\textsc{sg}.\Trs{} bread.\Dir{} child.\Ind{} eat-\Caus{}-\Imp{}//
\glft \trsl{You make the child eat the bread.} \lingcontext{overt subject in \Trs{}}//
\endgl
\xe

With an intransitive base, the overt causee itself takes direct case. The
following constructed pair holds that selection constant across the two
reduced aspectual values; the PFV exponent specifies a bounded directive:

\pex\label{ex:val-directive-caus-intr}
\a
\begingl
\gla Byl zemácrím.//
\glb child.\Dir{} run-\Caus{}-\Imp{}//
\glft \trsl{Make the child run.} \lingcontext{IPFV directive}//
\endgl
\a
\begingl
\gla Byl zemácrovím.//
\glb child.\Dir{} run-\Caus{}-\textsc{pfv.dom}-\Imp{}//
\glft \trsl{Make the child run.} \lingcontext{bounded PFV directive}//
\endgl
\xe

\noindent Without an overt causee the causative is ungrammatical in directive clauses:

\pex
\rec{trava pjáštrím} \trsl{intended: `Make someone eat the bread.'} \lingcontext{causative without causee}
\xe

The requirement for overt applied participants and causees is a constructional
restriction of directives. Since several indirect-case DPs may occur together,
a directive causative of an applicative can express the theme in direct case
and both beneficiary and causee in indirect case. When the base has no theme,
the highest eligible participant is direct: the applied participant takes
precedence over a causee. Morphological argument introduction consequently
remains distinct from the selection imposed by higher Mood.

Antipassive morphology is excluded from all directive moods by a
language-specific licensing restriction. Its availability in the indicative
ERG domain does not extend to directive PFV \ird{-ov-}; neither a silent
subject nor a reduced aspect inventory would independently explain this ban.
Commands with an unexpressed patient use the ordinary directive construction,
without overt antipassive morphology.

\section{Irrealis}\label{sec:irrealis}

Affirmative irrealis \ird{-ěv} has segmental content /ɛv/ and selects
the mixed consonantal grade together with local quantitative selection. Where the lexical root is exposed, a
licensed mid-vowel root selects raised-strong G3: \ird{zděk-}
\trsl{bring} gives \ird{zďíčev}. Root /ɛ/ becomes long /iː/, while
lexical /ɟ/, written \ird{dě} before /ɛ/, remains palatal and is written
\ird{ďí} before /iː/. Both properties are independent of the final velar
correspondence /k/ $\rightarrow$ /tʃ/.
Low-vowel \ird{pjašt-} instead supplies strong \ird{pjáščev},
with quantity on /a/ and a listed cluster alternant. Its consonantal
identity with the antipassive does not insert antipassive Voice into
the irrealis.

\pex
\a \ird{zděk-} + \ird{-ěv} $\rightarrow$ \ird{zďíčev}
\lingcontext{root raised-strong grade; independent mixed velar output}
\a \ird{pjašt-} + \ird{-ěv} $\rightarrow$ \ird{pjáščev}
\lingcontext{root quantity strength; listed modal cluster stem}
\xe

Intervening valency or orientation closes the root domain to this later
Mood selector. With compact Caus, \ird{pjáštr-} retains the G1
already supplied by Caus, but irrealis does not assign a new root grade:
\ird{pjáštrev} preserves /r/ and takes plain \ird{e}. Expanded
causatives can bear local strength on their own vowel; an already long
\ird{-ér-} retains that quantity, without raising the root through Caus.
The mixed consonantal map at a rhotic edge leaves /r/ unchanged.

Reduced perfective Asp supplies its own accessible nucleus. Affirmative
irrealis selects strong \ird{-óv-} from \ird{-ov-}, then applies
PAL to its final /v/, giving \ird{pjaštóvěv}. The lexical root
remains short and unmutated. The same local selection gives
\ird{pjáštróvěv} after compact Caus and
\ird{pjáštranóvěv} where causative-conditioned orientation
\ird{-an-} is present. The retained causative quantity and newly
selected Asp quantity belong to different domains, so their coexistence
does not license rootward reassignment or global shortening.

Negative irrealis \ird{-ir}, with segmental content /ɪr/, selects the
restricted velar-front consonantal grade without a new root quantitative grade. Thus
\ird{zděcir} retains the mid vowel, and \ird{pjaštir} retains
the short root and plain coronal. It also leaves reduced Asp
\ird{-ov-} short: \ird{pjaštovir},
\ird{pjáštrovir}, and \ird{pjáštranovir}. Previously selected
Caus or lexical quantity remains intact. Negative Mood realises
negation without \ird{za-}; it does not replace the independently
licensed inner stem.

\pex
\a \ird{pjaštir} \trsl{would/might not eat; irrealis non-eating}
\a \ird{zděcir} \trsl{would not bring}
\lingcontext{restricted k $\rightarrow$ c; no raised-strong root grade}
\xe

The affirmative irrealis is compatible with the converb-mood
constructions described in Chapter~\ref{ch:complex} and serves as the
main-clause mood in concessive constructions. Its distribution does
not alter the distinction between an exposed root and the edge of an
already assembled complex.
